\begin{table}[!htbp]
\centering
\footnotesize
\setlength{\tabcolsep}{3.5pt}
\caption*{\textbf{Table C10.} Out-of-zone tolerance measured from the feasibility boundary}
\begin{tabularx}{\linewidth}{lCCCC}
\toprule
Excess tolerance & Largest minimum band & Smallest main-band margin & All Seoul & Without the three districts \\
\midrule
\addlinespace[3pt]\multicolumn{5}{l}{\textit{2022}} \\
0.02 & 0.0306 & 0.0694 & [0.002, 0.198] & \textbf{[$-$0.013, 0.180]} \\
0.05 & 0.0306 & 0.0694 & \textbf{[$-$0.010, 0.208]} & \textbf{[$-$0.026, 0.190]} \\
0.10 & 0.0306 & 0.0694 & \textbf{[$-$0.031, 0.228]} & \textbf{[$-$0.049, 0.210]} \\
\addlinespace[3pt]\multicolumn{5}{l}{\textit{2024}} \\
0.02 & 0.0463 & 0.0537 & [0.028, 0.183] & [0.001, 0.161] \\
0.05 & 0.0463 & 0.0537 & [0.016, 0.193] & \textbf{[$-$0.012, 0.171]} \\
0.10 & 0.0463 & 0.0537 & \textbf{[$-$0.004, 0.211]} & \textbf{[$-$0.033, 0.190]} \\
\addlinespace[3pt]\multicolumn{5}{l}{\textit{2025}} \\
0.02 & 0.0606 & 0.0394 & \textbf{[$-$0.018, 0.177]} & \textbf{[$-$0.044, 0.156]} \\
0.05 & 0.0606 & 0.0394 & \textbf{[$-$0.034, 0.191]} & \textbf{[$-$0.061, 0.171]} \\
0.10 & 0.0606 & 0.0394 & \textbf{[$-$0.059, 0.214]} & \textbf{[$-$0.088, 0.195]} \\
\addlinespace[3pt]\multicolumn{5}{l}{\textit{2026}} \\
0.02 & 0.0732 & 0.0268 & [0.069, 0.141] & [0.044, 0.124] \\
0.05 & 0.0732 & 0.0268 & [0.057, 0.152] & [0.032, 0.136] \\
0.10 & 0.0732 & 0.0268 & [0.040, 0.170] & [0.012, 0.155] \\
\bottomrule
\end{tabularx}
\par\vspace{2pt}{\raggedright\footnotesize\textit{Notes.} Calibrated movement and residents' in-zone composition slack s=0.02 are held fixed. Each jurisdiction's minimum out-of-zone band is found by bisection (feasible upper endpoint; width <=0.00001). The main-band margin is 0.10 minus that minimum. Bounds use the jurisdiction-specific minimum plus the stated excess tolerance, capped at one. Thus these scenarios need not nest the constant 0.10 scenario. Compatibility with the counts is not validation of a restriction. Ranges concern the within-district roll-based slope; bold ranges include zero.\par}
\end{table}
